\section{Logical Framework of the Paper}
\label{sec:logical_framework}

This paper is structured to systematically address the challenges of NL-\nlwhat data preparation. The logical flow is as follows:

\begin{itemize}
    \item \textbf{Section 1: Introduction}
    \begin{itemize}
        \item Introduce the problem domain of Data Preparation (DP), highlighting its time-consuming nature.
        \item Propose Natural Language (NL) as an intuitive interface for specifying data transformations.
        \item Identify the core challenge: \textbf{Compositional Complexity}, which is broken down into:
        \begin{itemize}
            \item The search space explosion problem of transformation sequences.
            \item The supervision data scarcity problem for training models on full-chain tasks.
        \end{itemize}
        \item Review the limitations of existing approaches (Task-Specific Solutions and Guidance-Dependent Frameworks).
        \item Introduce our solution, \textbf{DeepPrep}, and its two core innovations designed to tackle the core challenge:
        \begin{itemize}
            \item Bi-Directional Full-Chain Supervision Synthesis.
            \item Tree-Structured Compositional Policy Learning.
        \end{itemize}
        \item Summarize the main contributions of the paper.
    \end{itemize}

    \item \textbf{Section 2: Preliminary}
    \begin{itemize}
        \item Formally define the task of Natural Language-\nlwhat Data Preparation, including the input, output, and transformation requirements.
        \item Define the \textbf{Logical Operation Space}, a comprehensive set of 31 standardized operators categorized for data preparation tasks.
        \item Provide a formal \textbf{Search Space Complexity Analysis} for traditional linear sequence approaches, mathematically demonstrating the intractability that motivates our work.
    \end{itemize}

    \item \textbf{Section 3: An Overview of \sys}
    \begin{itemize}
        \item Present the high-level architecture of the \sys framework.
        \item Introduce its three interconnected core components:
        \begin{itemize}
            \item \textbf{\promptframework:} The central reasoning engine that constructs a transformation tree.
            \item \textbf{\datasyn:} The data synthesis pipeline for generating training data.
            \item \textbf{\policylearn:} The training strategy for the orchestration model.
        \end{itemize}
    \end{itemize}

    \item \textbf{Section 4: Tree-based Operation Orchestration Framework (\promptframework)}
    \begin{itemize}
        \item Detail the core mechanism of the framework, which builds an operation tree through an iterative, multi-turn process.
        \item Explain its key capabilities that overcome the limitations of linear approaches:
        \begin{itemize}
            \item \textbf{Exploration:} The ability to pursue multiple transformation paths.
            \item \textbf{Rollback:} The ability to intelligently backtrack from errors.
        \end{itemize}
        \item Provide a theoretical analysis of how the tree-based structure achieves an \textbf{exponential reduction in search complexity} compared to linear methods.
    \end{itemize}

    \item \textbf{Section 5: Bi-Directional Full-Chain Supervision Synthesis (\datasyn)}
    \begin{itemize}
        \item Address the data scarcity problem by introducing a novel data synthesis pipeline.
        \item Describe the \textbf{Backward Data Synthesis} process, which uses an operation-guided inverse synthesis strategy to create realistic "dirty" data and corresponding cleaning sequences.
        \item Describe the \textbf{Forward Data Synthesis} process, which uses a multi-agent framework to translate SQL queries into our defined operation sequences.
        \item Explain how both directions are integrated to generate complete, end-to-end training examples with full-chain supervision.
    \end{itemize}

    \item \textbf{Section 6: Multi-Stage Policy Learning (\policylearn)}
    \begin{itemize}
        \item Propose a sophisticated training strategy to effectively teach the model complex reasoning.
        \item Detail the \textbf{Multi-Stage Cold Start} (SFT) process, which progressively trains the model from basic operator learning to complex multi-solution trajectory understanding.
        \item Introduce the \textbf{Multiturn RL with Hybrid Reward} phase, which refines the model's policy using:
        \begin{itemize}
            \item A hybrid reward function combining outcome-based, partial similarity, and LLM-as-Judge rewards to provide dense and meaningful feedback.
            \item A multiturn RL setup (GRPO) to learn long-horizon planning, exploration, and rollback.
        \end{itemize}
    \end{itemize}
\end{itemize}
